% Generated draft: requires mathematical and source review.
\documentclass[11pt]{article}
\usepackage[margin=25mm]{geometry}
\usepackage{fontspec}
\setmainfont{Latin Modern Roman}
\usepackage{amsmath,amssymb,mathtools,unicode-math}
\setmathfont{Latin Modern Math}
\usepackage{enumitem,needspace,xurl,hyperref}
\setlength{\parindent}{0pt}
\setlength{\parskip}{5pt}
\setlength{\emergencystretch}{3em}
\newcommand{\N}{\mathbb{N}}
\newcommand{\Z}{\mathbb{Z}}
\newcommand{\Q}{\mathbb{Q}}
\newcommand{\R}{\mathbb{R}}
\newcommand{\C}{\mathbb{C}}
\newcommand{\F}{\mathbb{F}}
\newcommand{\A}{\mathbb{A}}
\newcommand{\PP}{\mathbb{P}}
\newcommand{\E}{\mathbb{E}}
\newcommand{\eps}{\varepsilon}
\newcommand{\dd}{\,\mathrm{d}}
\begin{document}
\section*{UnsolvedMath ID: 30000263}
\textbf{Source problem code:} catalogue entry
\subsection{Definitions and mathematical statement}
For each positive integer $p$, let $a=(a_1,\ldots,a_p)\in(\R^3)^p$ satisfy $a_i\ne a_j$ whenever $i\ne j$. Write $a_i=(a_i^1,a_i^2,a_i^3)$ and define the real symmetric matrix $A(a)=(\alpha_{ij}(a))_{i,j=1}^p$ by
\[
 \alpha_{ii}=0,\qquad \alpha_{ij}=\frac1{|a_i-a_j|}\quad(i\ne j).
\]
All vector norms are Euclidean. The matrix derivative $\partial A/\partial a_i^k$ differentiates the entries with respect to coordinate $k\in\{1,2,3\}$ of point $a_i$, holding the other coordinates fixed. For $u=(u_1,\ldots,u_p)\in\R^p$, this gives
\[
 u^{\mathsf T}\frac{\partial A}{\partial a_i^k}u
 =-2u_i\sum_{j\ne i}\frac{u_j(a_i^k-a_j^k)}{|a_i-a_j|^3}.
\]

The conjecture is the following uniform inequality:
\[
 \forall p\in\N_{>0}\ \exists c(p)>0\ \forall a\in(\R^3)^p
 \text{ with pairwise distinct entries}\ \forall u\in\R^p,
\]
\[
 \max_{\substack{1\le i\le p\\1\le k\le3}}
 \left|u^{\mathsf T}\frac{\partial A(a)}{\partial a_i^k}u\right|
 +|A(a)u|^2
 \ \ge\ c(p)\sum_{\substack{1\le i,j\le p\\i\ne j}}
 \frac{u_i^2}{|a_i-a_j|^2}.
\]
The finite maximum is the source's $\sup_{i,k}$. The right-hand sum is over ordered pairs, exactly as written in the source. The constant may depend on $p$ but must be independent of both the configuration and $u$; no lower bound on separation or bound on diameter is added. The substantive cases have $p\ge2$; allowing $p=1$ gives the vacuous zero inequality.

The supplied literature assessment also mentions a proposed continuous limit, but no continuous-limit formulation is present in the source record. It is therefore recorded as missing context, not reconstructed as another conjecture.
\subsection{Short English statement}
The question asks for a positive constant depending only on the number of points that controls the weighted squared coordinates of every real vector by the interaction matrix and its coordinate derivatives, uniformly over all distinct configurations in three dimensions.
\subsection{Sources}
\begin{itemize}
\item Local UnsolvedMath registry, permanent ID 30000263, code catalogue entry, ``Discrete Interaction-Matrix Inequalities in Three Dimensions.'' Proposed year 2005; proposer unspecified.
\item \emph{Topological and Variational Methods for Differential Equations} (2006), Oberwolfach Reports, pp.\ 1601--1678, DOI 10.4171/owr/2005/29. Original problem report: \url{https://doi.org/10.4171/owr/2005/29}. The displayed report year 2006 differs from the DOI and proposed year 2005.
\item Supplied statement status corrected\_verified. The verification note says the official report was checked directly and the interaction matrix and full displayed conjectural inequality were transcribed. This verification claim is not independently checked here.
\item Literature triage checked 2026-08-21: open. It calls the item a conjectural multi-point inequality with a proposed continuous limit and says the supplied search found no published theorem resolving the exact stated inequality. Created 2026-08-25; updated 2026-09-16. Time-sensitive status and cited claims require independent verification.
\item Registry curation by Ulam AI / UnsolvedMath Contributors, CC BY 4.0; the underlying report retains its supplied provenance. Statement-only draft awaiting mathematical review.
\end{itemize}
\end{document}
